\documentclass[11pt,a4paper]{article}
\usepackage{common}
\setrunninghead{Microelectronics ECC-305 -- PBL Component I}
\captionsetup[figure]{name=Fig.}
\renewcommand{\thefigure}{R\arabic{figure}}
\hypersetup{pdftitle={ECC-305 PBL Component I -- Review of Break-before-Make CMOS Inverter},pdfauthor={}}

\begin{document}

\collegeheader

\begin{center}
  {\bfseries MICROELECTRONICS -- ECC-305 -- PBL COMPONENT -- I}\\[3pt]
  {\bfseries REVIEW OF RESEARCH PAPERS ON LOW-POWER TECHNOLOGY}\\[3pt]
  \textbf{Academic Year:} 2026--2027
\end{center}

\studentdetails
\evaluatorbox

\vspace{4pt}
\section*{\centering TITLE OF RESEARCH PAPER REVIEWED}
\begin{center}\textbf{Break-before-Make CMOS Inverter for Power-Efficient Delay Implementation}\end{center}

\section*{\centering DETAILS OF RESEARCH PAPER}
\textbf{Title:} Break-before-Make CMOS Inverter for Power-Efficient Delay Implementation

\textbf{Author(s):} Janez Puhan, Du\v{s}an Rai\v{c}, Tadej Tuma, and \'Arp\'ad B\H{u}rmen (Faculty of Electrical Engineering, University of Ljubljana, Slovenia)

\textbf{Journal/Conference:} The Scientific World Journal (Hindawi Publishing Corporation), Volume 2014

\textbf{Year of Publication:} 2014 (received 4 June 2014; accepted 3 October 2014; published 26 November 2014)

\textbf{DOI / Paper Identification:} Article ID 349131, 11 pages; DOI 10.1155/2014/349131; PMC4265520. Topic assigned: Low-Power Technology, Break-Before-Make (BBM) CMOS Inverter.

\newpage
\section*{\centering REVIEW SUBMISSION}
The review includes the following parts.

% ------------------------------------------------------------------
\itemhead{Background and motivation of the research}
Digital systems often need a small, low-precision delay (for example, to align or strobe signals). The usual solution is a chain of inverters. The paper observes that an inverter chain is power-efficient only for small delays: every inverter in the chain drains additional parasitic energy approximately equal to the dynamic energy needed to change its input, so long delays become power-hungry. Direct-path (short-circuit) current is a well-known source of internal dynamic power; in well-designed circuits it is estimated at under \SI{20}{\percent} of dynamic dissipation (paper ref.~[1], Veendrick 1984), but it can become prohibitive with significant capacitive loads. The motivation is therefore to obtain a required delay with the minimum supply charge (energy) by suppressing direct-path current.

\itemhead{Problem addressed by the researchers}
How can a required delay \Tp{} be implemented at the lowest possible charge $Q$ at gate level, where loads are small? Earlier remedies drive the NMOS and PMOS gates of an inverter with separate, time-skewed signals (paper refs.~[2]--[5]), but all of them insert additional driver stages ahead of the split inverter; that overhead in area and power is justified only for large capacitive loads, not for gate-level logic. Other gate-level techniques ([6], [7]) address static leakage in nanometre technologies, carry an area overhead, and do not improve internal dynamic power. The paper targets this gap: internal dynamic power of small-load gate-level delay elements.

\itemhead{Low-power technology/technique discussed}
Break-before-make (BBM) switching, also called dual-ramp (time-skewed) drive, eliminates direct-path (short-circuit) current. The PMOS and NMOS switching transistors of an inverter are never ON at the same time: during an ``isolation time'' both are off and the output node is high-impedance, holding its previous logic state on the load capacitance. The technique is dynamic-power oriented (short-circuit component), not a leakage or supply-scaling technique.

\itemhead{Existing approach}
A standard two-stage CMOS buffer or serial inverters (paper Fig.~2). The charge per cycle is the sum of a dynamic part in the first stage and a direct-path part $Q_s$ in the second stage, whose input rise/fall times $t_r$, $t_f$ (Eq.~5) are non-zero:
\begin{equation}\tag{4}
  Q_d = (C_{LP}+C_{LN})\,V_{DD}.
\end{equation}
For a symmetric cell ($V_{TN}=-V_{TP}=V_T$, $\beta_N=\beta_P$) the paper obtains
\begin{equation}\tag{9}
  Q = Q_d + Q_s = \left[\frac{2\beta_1 V_{DD}^{2}}{k_2} + \frac{\beta_2}{6V_{DD}}\,(V_{DD}-2V_T)^{3}\right] T_{p,\mathrm{CMOS}},
\end{equation}
\begin{equation}\tag{10}
  T_{p,\mathrm{CMOS}} = \frac{k_2\,(C_{LP}+C_{LN})}{2\beta_1 V_{DD}}.
\end{equation}
Both static (direct-path) and dynamic charge grow linearly with the CMOS propagation delay. Here $\beta = (\mu\varepsilon/t_{ox})(w/l)$ and $k_2$ is a constant; the analysis uses the simple model of Eq.~3 (cut-off / triode / saturation square law).

\begin{figure}[H]
  \centering
  \includegraphics[width=0.7\textwidth]{paper_fig9.png}
  \caption{Standard (a) and proposed (b) delay cell with unit input/output inverters. Reproduced from Puhan et al.\ (2014), Fig.~9, CC~BY licence.}
\end{figure}

\itemhead{Proposed approach}
A bidirectional delay element is connected in series between the PMOS and NMOS drains of the first inverter, producing two time-skewed outputs ($p_o$, $n_o$). The next stage is an inverter with split inputs ($p_i$, $n_i$) whose outputs are merged (Fig.~1(c) of the paper). Stages cascade without accumulating isolation time, and the overhead is only the delay element, with no extra driver stages.
\begin{itemize}
  \item \textbf{$R_d$ version.} A resistor $R_d$ is the simplest delay element; its $RC$ product with the next-stage gate capacitance sets the isolation time (Eq.~14: $t_{if}=k_8R_dC_{LP}$, $t_{ir}=k_8R_dC_{LN}$). Polysilicon $R_d$ is area-hungry and poorly controlled.
  \item \textbf{Transistor versions (paper Fig.~6).} A diode- or triode-connected transistor reduces the output swing by a threshold voltage, needs $V_{DD}>5V_{Td}$ and gives long information transients, so reduced-swing topologies are judged inappropriate.
  \item \textbf{Full-swing versions (paper Figs.~7, 8).} Level-restoration transistors (controlled by the next-stage output) are added in parallel with the delay transistor. Four two-transistor topologies are compared: PtNt, PfNf, PtNf and PfNt. PtNf (a PMOS delay transistor in triode plus an NMOS level-restoration transistor) is reported as the most efficient.
\end{itemize}

\begin{figure}[H]
  \centering
  \includegraphics[height=6.5cm]{paper_fig8_PtNf.png}
  \caption{PtNf full-swing delay implementation. Reproduced from Puhan et al.\ (2014), Fig.~8, CC~BY licence.}
\end{figure}

\itemhead{Important circuit/device concepts}
\begin{figure}[H]
  \centering
  \includegraphics[width=0.95\textwidth]{paper_fig1.png}
  \caption{BBM structures with time-domain responses: (a) CMOS-to-BBM converter, (b) BBM inverter, (c) BBM-to-CMOS converter. Reproduced from Puhan et al.\ (2014), Fig.~1, CC~BY licence.}
\end{figure}

\begin{itemize}
  \item \textbf{Isolation versus information slopes.} The first edge of a skewed pair (falling edge on the NMOS gate, rising edge on the PMOS gate) is the isolation slope; the second is the information slope. The interval between them is the isolation time $t_i$, during which the output is high-impedance and no direct path exists. If $t_i\to 0$ the circuit becomes an ordinary CMOS inverter. Isolation times do not accumulate through a cascade.
  \item \textbf{Redefined delays.} $t_{dr}$, $t_{df}$ are measured from the input information slope to the output isolation slope; $t_{ir}$, $t_{if}$ lie between isolation and information slopes. The propagation delay (paper Eqs.~1 and 2, symmetric case) is
\end{itemize}
\begin{equation}\tag{1, 2}
  T_p = \frac{t_{ir}+t_{df}+t_{if}+t_{dr}}{2} = t_d + t_i, \qquad T_p = k_1\,T_{p,\mathrm{CMOS}} + t_i,\quad k_1<1.
\end{equation}
\begin{itemize}
  \item \textbf{Charge of the $R_d$ cell.} With $V_{TN}=-V_{TP}$, a symmetric first stage and high $R_d$, the static charge is $Q_s\approx 0$ and (Eqs.~13--16) $t_{df}=k_6C_{LN}/(\beta_{N1}V_{DD})$, $t_{dr}=k_7C_{LP}/(\beta_{P1}V_{DD})$ and
\end{itemize}
\begin{equation}\tag{16}
  Q = \frac{2\beta_1 V_{DD}^{2}}{k_6 + k_8R_d\beta_1V_{DD}}\;T_p .
\end{equation}
\begin{itemize}
  \item \textbf{Condition for valid BBM operation (Eqs.~17--18).} The isolation slope must end before the information slope begins, i.e.
\end{itemize}
\begin{equation}\tag{17, 18}
  R_d > \max\!\left(\frac{k_2C_{LN}}{k_{11}C_{LP}\beta_{N1}V_{DD}},\;\frac{k_3C_{LP}}{k_{11}C_{LN}\beta_{P1}V_{DD}}\right).
\end{equation}
\begin{itemize}
  \item \textbf{Figures of merit.} PDP (energy per switching event) and EDP $=\mathrm{PDP}\cdot T_p$; for a delay line the goal is minimal charge $Q$ at a required \Tp{}, since \VDD{} is fixed.
  \item \textbf{Level restoration.} A reduced-swing node (a $V_{Td}$ drop) is restored by $N_f$/$P_f$ devices driven from the next-stage output (feedback).
\end{itemize}

\itemhead{Simulation/experimental methodology}
\begin{itemize}
  \item \textbf{Simulator and model.} Synopsys HSPICE with TSMC \SI{0.18}{\micro\meter} / \SI{1.8}{\volt} process parameters (proprietary; I did not reproduce it).
  \item \textbf{Circuits compared.} The standard cell of Fig.~9(a) (four inverters) against four BBM delay cells (PtNt, PtNf, PfNt, PfNf) of Fig.~9(b), each with identical standard input and output unit inverters so that fan-in and fan-out are equal; both unit inverters are counted as part of the cell. An 8-dimensional (standard) or 12-dimensional (BBM) sizing problem was solved for each delay target from \SI{100}{\pico\second} to \SI{5}{\nano\second}.
  \item \textbf{Optimiser.} Parallel SADE global optimisation (paper ref.~[12]) driven by PyOPUS, on 100 computing nodes (25 Intel Core i5 \SI{2.66}{\giga\hertz} processors, 4 nodes each); every run was repeated from several starting points until the global minimum was confirmed.
  \item \textbf{Constraints.} (i) Limits on over/undershoot so the circuit does not rely on Miller (parasitic) capacitance; (ii) steady state reached after $t_{\max}$; (iii) noise-margin constraints from p- and n-substrate disturbances. Without them the optimiser returns ``degenerate'' circuits with very long transients.
  \item \textbf{Parasitics.} Layout was not in the optimisation loop; instead the layout rules $A_d=A_s=\SI{0.8}{\micro\meter}\times w$ and $P_d=P_s=\SI{1.5}{\micro\meter}+w$ were applied for source/drain parasitics.
  \item \textbf{Metric.} Charge consumption (fAs) for one rising and one falling slope; Table~1 below.
\end{itemize}

\itemhead{Key results reported in the paper}
\textbf{Table 1} of the paper gives the charge consumption of the standard (std.) and BBM delay cells of Fig.~9, measured for one rising and one falling slope, in fAs (femto-ampere-seconds); \% is the percentage of the standard cell's consumption. Values are transcribed from the paper; ``*'' marks a delay target that was not achieved. Conditions: HSPICE, TSMC \SI{0.18}{\micro\meter} / \SI{1.8}{\volt}, each cell individually sized for the target \Tp{}.

\begin{table}[H]
  \centering\small
  \caption*{\textbf{Paper Table 1.} Charge consumption (fAs) of the standard and BBM delay cells.}
  \begin{tabular}{@{}r r r r r r r r r r@{}}
    \toprule
    $T_p$ [ns] & std. & PtNt & \% & PtNf & \% & PfNt & \% & PfNf & \%\\
    \midrule
    0.1  & 34*  & 45*  & --- & 78* & --- & 44* & --- & 97* & ---\\
    0.2  & 26   & 30   & 115 & 87* & --- & 30  & 115 & 62* & ---\\
    0.35 & 32   & 34   & 106 & 34  & 106 & 35  & 109 & 42  & 131\\
    0.5  & 39   & 38   & 97  & 35  & 90  & 39  & 100 & 38  & 97\\
    0.75 & 65   & 45   & 69  & 42  & 65  & 46  & 71  & 46  & 71\\
    1.0  & 77   & 58   & 75  & 49  & 64  & 58  & 75  & 55  & 71\\
    2.0  & 113  & 86   & 76  & 67  & 59  & 86  & 76  & 93  & 82\\
    3.5  & 153  & 118  & 77  & 89  & 58  & 118 & 77  & 130 & 85\\
    5.0  & 186  & 142  & 76  & 108 & 58  & 143 & 77  & 177 & 95\\
    \bottomrule
  \end{tabular}
\end{table}

\begin{itemize}
  \item Charge grows approximately linearly with delay above about \SI{1}{\nano\second} (as Eqs.~9 and 16 predict). PtNf is the most efficient topology; the charge saving against the standard cell is slightly above \SI{40}{\percent} (\SI{58}{\percent} of std.\ at \SI{3.5}{\nano\second} and \SI{5}{\nano\second}). The abstract states that the energy (charge) per delay is reduced ``up to 40\,\%''.
  \item With the exception of very short delays the BBM structure improves the power budget: at \SIrange{0.1}{0.35}{\nano\second} the BBM cells are no better (PtNf \SI{106}{\percent} at \SI{0.35}{\nano\second}) or do not reach the target delay.
  \item Third-stage drain currents for $T_p=\SI{750}{\pico\second}$ (paper Fig.~11): the large direct-path current of the standard cell is almost completely eliminated in PtNf. The sizing in the Fig.~11 caption is: \emph{std.}\ $w_{P1}=\SI{220}{\nano\meter}$, $l_{P1}=\SI{370}{\nano\meter}$, $w_{N1}=\SI{220}{\nano\meter}$, $l_{N1}=\SI{880}{\nano\meter}$, $w_{P2}=\SI{880}{\nano\meter}$, $l_{P2}=\SI{730}{\nano\meter}$, $w_{N2}=\SI{1.05}{\micro\meter}$, $l_{N2}=\SI{1.32}{\micro\meter}$; \emph{PtNf:} $w_{P1}=\SI{220}{\nano\meter}$ / $l_{P1}=\SI{740}{\nano\meter}$, $w_{N1}=\SI{390}{\nano\meter}$ / $l_{N1}=\SI{180}{\nano\meter}$, $w_{Pt}=\SI{220}{\nano\meter}$ / $l_{Pt}=\SI{350}{\nano\meter}$, $w_{Nf}=\SI{220}{\nano\meter}$ / $l_{Nf}=\SI{950}{\nano\meter}$, $w_{P2}=\SI{220}{\nano\meter}$ / $l_{P2}=\SI{180}{\nano\meter}$, $w_{N2}=\SI{220}{\nano\meter}$ / $l_{N2}=\SI{220}{\nano\meter}$. The bulk of the delay comes from the third-stage gate capacitances.
\end{itemize}

\textbf{Application results.} The paper presents a delay line replacing inverter chains (Fig.~12) and a leading-edge pulse generator (Fig.~13). Paper Table~2 gives the charge (fAs) per input impulse, standard versus BBM-based; ``*'' marks a pulse-width target that was not achieved.

\begin{table}[H]
  \centering\small
  \caption*{\textbf{Paper Table 2.} Charge (fAs) per input impulse, standard versus BBM-based pulse generator.}
  \begin{tabular}{@{}r r r r r r r r@{}}
    \toprule
    No.\ of inv. & Width [ps] & Std total & Std delay & BBM total & \% & BBM delay & \%\\
    \midrule
    3  & 200 & 171 & 56  & 143* & --- & 17* & ---\\
    5  & 330 & 210 & 93  & 142  & 68  & 16  & 17\\
    7  & 450 & 247 & 130 & 147  & 60  & 16  & 12\\
    9  & 570 & 284 & 167 & 153  & 54  & 16  & 10\\
    11 & 690 & 321 & 204 & 161  & 50  & 17  & 8\\
    \bottomrule
  \end{tabular}
\end{table}

The paper concludes that the BBM-based leading-edge pulse generator saves up to \SI{50}{\percent} of the total switching energy. The delay-line charge of the BBM version is almost constant (\SIrange{16}{17}{fAs}) because the number of transistors is constant, whereas the standard delay line grows linearly with the number of inverters. Layout parasitics of the cell nodes were not included.

\itemhead{Advantages of the proposed approach}
\begin{itemize}
  \item Direct-path current is practically eliminated in the driven stage, since the PMOS and NMOS are never ON simultaneously (paper Fig.~11).
  \item Very low component overhead: no additional driver stages; a two-transistor delay element suffices (a six-transistor non-inverting delay cell), so it is usable with the small loads of gate-level design.
  \item Charge saving of slightly above \SI{40}{\percent} at longer delays (PtNf) and up to \SI{50}{\percent} of total energy in the pulse generator.
  \item Delay and short-circuit elimination are obtained in one stage; two output signals give compact functional solutions (for example, an edge detector).
  \item Isolation times do not accumulate along a cascade, so BBM stages can be chained.
\end{itemize}

\itemhead{Limitations of the proposed approach}
\begin{itemize}
  \item Not beneficial for very short delays ($\leq\SI{0.35}{\nano\second}$: PtNf \SI{106}{\percent} at \SI{0.35}{\nano\second}; target not reached at \SIrange{0.1}{0.2}{\nano\second}). The authors say so themselves: for logic gates the extra delay and double wiring hardly justify the energy saving. It fits circuits that use the delay (edge-triggered storage, clock distribution).
  \item Requires dual (two-wire) signals between stages, which costs routing area.
  \item Reduced-swing transistor topologies need $V_{DD}>5V_{Td}$ and have long information transients; full swing requires extra level-restoration transistors with feedback from the next stage.
  \item Optimiser results depend on parasitics unless extra constraints are imposed; the authors could not include layout extraction in the optimisation loop, so layout-induced node capacitances are not accounted for.
  \item Only simulation results are presented (HSPICE, one proprietary \SI{0.18}{\micro\meter} process); no fabricated-silicon measurements, PVT/Monte-Carlo corners or leakage analysis are reported. The delay is ``low-precision'' and sensitive to the poorly controlled resistance of the delay element.
\end{itemize}

\itemhead{Scope for further improvement}
\begin{itemize}
  \item Include layout extraction and parasitics in the optimisation loop (named by the authors as limited only by computing power).
  \item Verify with process/voltage/temperature corners and Monte-Carlo mismatch, and with modern FinFET/sub-\SI{100}{\nano\meter} nodes where leakage dominates.
  \item Obtain better delay precision using the biased-NMOS (paper Fig.~14) or dynamically biased (paper Fig.~15) cells, whose DC bias power is reduced by dynamic biasing.
  \item Extend the concept to logic gates (NAND/NOR) in a way that reduces the double-wiring overhead, and to clock distribution and edge-triggered storage elements.
  \item My own idea: pair BBM with a lower \VDD{} or power gating, since BBM does nothing about leakage.
\end{itemize}

\itemhead{Relevance to Microelectronics}
The paper touches most of what the course covers: (i) the three components of CMOS power, dynamic ($CV_{DD}^2f$), short-circuit and leakage, and techniques to reduce the short-circuit term; (ii) square-law MOSFET operation regions (Eq.~3) and $\beta$, $V_T$, body effect and channel-length effects; (iii) propagation delay, rise/fall times, PDP and EDP; (iv) transistor sizing trade-offs (long-channel devices create delay cheaply but cost area); (v) level restoration and voltage swing; (vi) simulation-driven design with SPICE and optimisation; and (vii) practical circuits such as delay lines, pulse generators and edge detectors used in dynamic logic and flip-flops.

\itemhead{Technical observations and conclusions}
\textbf{Part A: my own simulation (kept separate from the paper's results).} The paper's TSMC \SI{0.18}{\micro\meter} HSPICE model is proprietary, so I rebuilt the Fig.~9 cells in ngspice-42 with a Level-1 model. Its main parameters come from the Fig.~4 caption of the paper: $V_{TN}=\SI{0.5}{\volt}$, $V_{TP}=\SI{-0.5}{\volt}$, $K_{PN}=\SI{350}{\micro\ampere\per\volt\squared}$, $K_{PP}=\SI{90}{\micro\ampere\per\volt\squared}$ and $C_{ox}=\SI{8.6}{\femto\farad\per\micro\meter\squared}$. Everything else is my own assumption: $\lambda$, $\gamma$, $\phi$, the overlap and junction capacitances, the unit inverters ($W_p/W_n=\SI{860}{\nano\meter}/\SI{220}{\nano\meter}$, $L=\SI{180}{\nano\meter}$) and a \SI{2}{\femto\farad} load. Supply \SI{1.8}{\volt}, \SI{27}{\celsius}, \SI{50}{\mega\hertz} input with \SI{50}{\pico\second} edges. $Q$ is the charge the whole cell draws from the supply over one input cycle, averaged over four steady-state cycles. Netlists and details are in the ECC-353 report.

\begin{table}[H]
  \centering\small
  \caption*{\textbf{Own ngspice results} (substitute Level-1 model).}
  \begin{tabular}{@{}p{7.2cm} r r r@{}}
    \toprule
    Case & $T_p$ (ps) & $Q$ (fC/cycle) & BBM charge saving\\
    \midrule
    Paper's \SI{750}{\pico\second} sizing, both cells unchanged, std.\ & 507 & 67.3 & ---\\
    Paper's \SI{750}{\pico\second} sizing, both cells unchanged, BBM PtNf & 563 & 40.3 & \SI{40.2}{\percent} (delays \emph{not} matched)\\
    Re-optimised for $T_p\approx\SI{0.5}{\nano\second}$, std.\ / BBM & 495 / 505 & 44.1 / 38.3 & \SI{13.1}{\percent}\\
    Re-optimised for $T_p\approx\SI{1.0}{\nano\second}$, std.\ / BBM & 1016 / 1006 & 64.8 / 49.8 & \SI{23.2}{\percent}\\
    \bottomrule
  \end{tabular}
\end{table}

\textbf{Part B: how this compares with the paper.} Table~1 of the paper puts PtNf at \SI{90}{\percent} of the standard cell's charge at \SI{0.5}{\nano\second} (a \SI{10}{\percent} saving) and \SI{64}{\percent} at \SI{1.0}{\nano\second} (a \SI{36}{\percent} saving). With matched delays my model gives \SI{13.1}{\percent} and \SI{23.2}{\percent}. The shape is the same: the benefit grows with delay and is small near \SI{0.5}{\nano\second}. The numbers are not. My saving at \SI{1}{\nano\second} is clearly below \SI{36}{\percent}, and the absolute charges cannot be compared because the model, the unit inverters and the load all differ. My optimiser is also tiny next to the paper's 100-node SADE run, so neither of my sizings is likely to be a true optimum.

\begin{figure}[H]
  \centering
  \includegraphics[width=0.5\textwidth]{fig_wave_bbm_skew.png}
  \caption{My ngspice result: skewed pair A (PMOS gate) / B (NMOS gate), isolation edge first, information edge second (paper's \SI{750}{\pico\second} PtNf sizing, substitute model).}
\end{figure}

\textbf{Part C: what I noticed.} The skew is easy to see in the waveforms. With the paper's \SI{750}{\pico\second} sizing the isolation edge (A rising or B falling) arrives \SIrange{155}{258}{\pico\second} before the information edge, and the two conducting windows never overlap. What surprised me was the crowbar charge. In the conventional cell's third stage it is only about \SI{0.7}{\femto\coulomb} per cycle (a rough $\min(i_P,i_N)$ estimate) out of the \SI{15.4}{\femto\coulomb} that stage draws, so I cannot say the saving comes mainly from removing short-circuit current. A good part probably comes from BBM letting $P_2$ and $N_2$ stay small. Node A also overshoots to \SI{2.04}{\volt} through Miller coupling, which is exactly what the paper's constraints are meant to suppress. Finally, taking the delay element out is not a fair ablation: $T_p$ drops from \SI{563}{\pico\second} to \SI{155}{\pico\second}, which only shows that the $P_t$ resistance is what makes the delay.

\textbf{Part D: conclusions.} BBM switching does lower the charge needed for a given delay. The paper reports slightly more than \SI{40}{\percent} at \SI{3.5}{\nano\second} and above, and nothing gained below about \SI{0.35}{\nano\second}. My substitute-model runs show the same trend, with \SI{23}{\percent} saving at about \SI{1}{\nano\second} and \SI{13}{\percent} at about \SI{0.5}{\nano\second}. For ordinary logic gates it is a poor deal, since the delay is not wanted and the second signal wire costs area. My results depend on the Level-1 model (no velocity saturation, no leakage), the small load and a short optimisation, so I treat them as a trend check, not a reproduction.

% ------------------------------------------------------------------
\begin{thebibliography}{9}
  \bibitem{puhan} J.~Puhan, D.~Rai\v{c}, T.~Tuma, and \'A.~B\H{u}rmen, ``Break-before-Make CMOS Inverter for Power-Efficient Delay Implementation,'' \emph{The Scientific World Journal}, vol.~2014, Article ID 349131, 11~pp., 2014. doi:\href{https://doi.org/10.1155/2014/349131}{10.1155/2014/349131}. \url{https://pmc.ncbi.nlm.nih.gov/articles/PMC4265520/} (primary source; CC~BY).
  \bibitem{veendrick} H.~J.~M. Veendrick, ``Short-circuit dissipation of static CMOS circuitry and its impact on the design of buffer circuits,'' \emph{IEEE J.\ Solid-State Circuits}, vol.~SC-19, no.~4, pp.~468--473, 1984 (cited in [1] as its ref.~[1]; not independently consulted).
  \bibitem{weste} N.~H.~E. Weste and D.~M. Harris, \emph{Principles of CMOS VLSI Design: A Circuits and Systems Perspective}, 4th ed., Addison-Wesley, 2011 (cited in [1] for the square-law model; not independently consulted).
  \bibitem{ngspice} ngspice, open-source mixed-level/mixed-signal circuit simulator, version~42, \url{https://ngspice.sourceforge.io/} (used for my simulations).
\end{thebibliography}

\end{document}
